% Generated draft: requires mathematical and source review.
\documentclass[11pt]{article}
\usepackage[margin=25mm]{geometry}
\usepackage{fontspec}
\setmainfont{Latin Modern Roman}
\usepackage{amsmath,amssymb,mathtools,unicode-math}
\setmathfont{Latin Modern Math}
\usepackage{enumitem,needspace,xurl,hyperref}
\setlength{\parindent}{0pt}
\setlength{\parskip}{5pt}
\setlength{\emergencystretch}{3em}
\newcommand{\N}{\mathbb{N}}
\newcommand{\Z}{\mathbb{Z}}
\newcommand{\Q}{\mathbb{Q}}
\newcommand{\R}{\mathbb{R}}
\newcommand{\C}{\mathbb{C}}
\newcommand{\F}{\mathbb{F}}
\newcommand{\A}{\mathbb{A}}
\newcommand{\PP}{\mathbb{P}}
\newcommand{\E}{\mathbb{E}}
\newcommand{\eps}{\varepsilon}
\newcommand{\dd}{\,\mathrm{d}}
\begin{document}
\section*{UnsolvedMath ID: 30006321}
\textbf{Source problem code:} catalogue entry
\subsection{Definitions and mathematical statement}
Let \(G\) and \(G^\vee\) denote the group and its Langlands dual intended in the supplied source. Let \(X\) denote the geometric datum defining the basic object \(\delta_X\), and let \(M^\vee\) denote its proposed dual geometric object, equipped with the relevant \(G^\vee\)-action and grading. The source does not specify the admissible class of these data, the base field, or the sheaf coefficients.

Write \(\mathcal C_X\) for the intended category containing \(\delta_X\), and \(\Phi\) for the applicable derived geometric Satake equivalence. Here \(\mathcal C_X\) is a placeholder for a category whose precise identification is unresolved below. Write \((-)^{\mathrm{sh}}\) for the source's shearing operation, which reinterprets grading in cohomological terms; its degree conventions are not supplied. The notation \(\mathcal O(M^\vee)\) denotes the structure sheaf or function algebra regarded as an object of the intended equivariant quasi-coherent category, whose precise realization also requires specification.

The intended conjectural assertion is that, for every datum in the source's unspecified range of applicability,
\[
\Phi(\delta_X)\simeq \mathcal O(M^\vee)^{\mathrm{sh}},
\]
and that the resulting identification of endomorphism algebras has the form
\[
\operatorname{End}_{\mathcal C_X}(\delta_X)
\simeq
\bigl(\mathcal O(M^\vee)^{G^\vee}\bigr)^{\mathrm{sh}}.
\]
These formulas retain the source's notation: whether endomorphisms mean a derived endomorphism object, a graded algebra, or degree-zero endomorphisms, and whether invariants are ordinary or derived, is not determined by the record.

The extracted passage places \(\delta_X\) in
\[
\operatorname{Shv}_{G(D)}\bigl(G(D^\circ)/G(D)\bigr),
\]
where \(D\) denotes a formal disk and \(D^\circ\) its punctured disk, and places the proposed image in an equivariant quasi-coherent category over \(\check{\mathfrak g}^{*}\), with \(\check{\mathfrak g}=\operatorname{Lie}(G^\vee)\) and an additional \(\mathbb G_m\)-grading. Elsewhere the same extraction expresses the endomorphisms using \(\operatorname{Shv}_{G_{\mathcal O}}(X_F)\) and \(\operatorname{Shv}_{G_{\mathcal O}}(X_{\mathcal O})\), with objects \(\delta_X\) and \(\omega_X\), and describes the answer as equivariant Borel--Moore homology of \(X_{\mathcal O}\), after shearing. The definitions of these spaces and objects, and the identifications connecting these categories, are absent. Consequently, the displayed conjecture is a conditional reconstruction of the intended target, not a fully specified universally quantified statement.
\subsection{Short English statement}
The proposed compatibility says that derived geometric Satake sends the basic object \(\delta_X\) to the sheared structure sheaf of the proposed dual space \(M^\vee\), with its endomorphisms corresponding to sheared invariant functions. The admissible geometric data, categories, shearing conventions, and meanings of endomorphisms and invariants require clarification. The supplied record labels the problem open; this draft does not independently verify that status.
\subsection{Sources}
\begin{itemize}
\item Local UnsolvedMath registry record supplied with this task: permanent ID \(30006321\), title ``Derived Geometric Satake for the Basic Object,'' exact source problem code \texttt{catalogue entry}, in the Oberwolfach Reports Open Problems collection. The record assigns the year 2025 and names no individual proposer.
\item \emph{Arbeitsgemeinschaft: Relative Langlands Duality} (2025), Oberwolfach Reports, pp.\ 811--882, DOI 10.4171/owr/2025/18. Supplied source report and literature link: \url{https://doi.org/10.4171/owr/2025/18}. No narrower location for the conjecture is supplied.
\item \emph{Relative Langlands Duality}, supplied literature link: \url{https://arxiv.org/abs/2409.04677}. The supplied assessment refers to this as the BZSV paper; no author list or conjecture number is supplied.
\item Literature assessment dated \texttt{2026-08-22}: the supplied curation labels the problem open and states that the BZSV paper and the 2025 report present the basic-object assertion as a conjecture, provide examples, and report no general proof. This is supplied dated curation, not an independently verified literature assessment. The background explicitly requests independent verification of time-sensitive status and cited claims.
\item The registry labels its statement \texttt{corrected\_verified} and reports that the official Oberwolfach report was checked and notation repaired. These are supplied provenance claims; no source visit or independent mathematical verification is asserted here.
\end{itemize}
\end{document}
